En este reporte se presenta el estudio experimental
del rendimiento práctico de algoritmos de ordenamiento
de arreglos unidimensionales (Merge Sort $\Theta(n \log n)$, 
Quick Sort $O(n^2)$, Patience Sort $O(n^2)$ y sort $\Theta(n \log n)$)
y de algoritmos de multiplicación de matrices cuadradas 
(Naive $\Theta(n^3)$ y Strassen $\Theta(n^{2.81})$).

\vspace{0.4cm}

Esto se realizo mediante la implementación de cada algoritmo
en C++ y la ejecución de pruebas controladas para medir su
rendimiento en términos de tiempo de ejecución y consumo máximo
de memoria RAM (\textit{Peak RSS}), el objetivo principal de 
tomar estas mediciones es contrastar las cotas asintóticas
teóricas con las complejidades observadas en la ejecucion de 
distintos casos, considerando factores como la arquitectura
de memoria, la profundidad de recursión y las condiciones 
especificas de los archivos de entrada.